\section{El desorden provocado por la avalancha de capital}

La sección anterior trató la productividad; esta trata el desorden en la industria. Después de ChatGPT, la inteligencia corporizada pasó de ser un concepto minoritario a convertirse en la favorita del capital, y el financiamiento, las valuaciones y los demo se dispararon. El núcleo de la crítica de Wang He (王鹤) es este: el mayor desorden de la industria consiste en no poder distinguir quién crea productividad y quién solo cuenta una historia. Más grave aún, algunos equipos no aclaran en sus presentaciones que usan teleoperación, e incluso hacen pasar la teleoperación por capacidad autónoma, lo que destruye directamente la confianza del público y de los clientes en la industria.

Este tipo de problema no es solo ético; también se convierte en un problema técnico. Las presentaciones engañosas desvían el reward model de la industria: los equipos optimizan los demo para conseguir financiamiento y difusión, en lugar de optimizar el sistema para la tasa de éxito en escenarios reales, la reducción de costos y el retorno de datos de las unidades entregadas. A largo plazo, el capital castigará a toda la industria, y los equipos que de verdad generan productividad también se verán arrastrados por el ruido.

\begin{figure}[H]
\centering
\includegraphics[width=0.92\textwidth]{capital-chaos-reward.png}
\caption{El desorden después de la avalancha de capital: quién crea productividad y quién solo cuenta una historia. Figura conceptual propia, elaborada a partir de la conversación de 02:13:51--02:25:25 y de la descripción del video.}
\end{figure}

\begin{knowledgebox}{Lectura de la figura: un reward equivocado produce conductas equivocadas}
En el centro de la figura está el reward. Si el reward proviene del monto de financiamiento, la valuación y el tráfico, los equipos optimizarán naturalmente la historia y la demostración; solo si el reward proviene de presentaciones veraces, de la recompra de los clientes y de aplicaciones de diez mil unidades, los equipos optimizarán la productividad. Al leer la figura, conviene considerar la «presentación veraz» como infraestructura de la industria: sin presentaciones confiables, el mercado no puede juzgar el progreso.
\end{knowledgebox}

\subsection{El reloj de validación: diez mil unidades en cinco años}

Wang He plantea un juicio contundente: en un plazo de cinco años, el campo debe contar con aplicaciones de más de diez mil unidades; si no lo logra, el campo de la inteligencia corporizada quedará refutado, o al menos lo quedará su narrativa actual. Esta afirmación no es necesariamente una predicción precisa, pero ofrece un criterio de aceptación para la industria. Las diez mil unidades no buscan una cifra vistosa, sino demostrar que el hardware, el software, la entrega, el mantenimiento, el retorno de datos y el valor para el cliente ya superaron la etapa de laboratorio.

\begin{figure}[H]
\centering
\includegraphics[width=0.96\textwidth]{robotics-validation-clock.png}
\caption{Reloj de validación de cinco años: se necesitan aplicaciones de más de diez mil unidades en cinco años; de lo contrario, la narrativa del campo quedará refutada. Figura conceptual propia, elaborada a partir de la conversación de 02:13:51--02:25:25 y de la descripción del video.}
\end{figure}

\begin{importantbox}{Principio de presentación veraz}
Se puede usar teleoperación, pero hay que decirlo; se pueden mostrar productos inacabados, pero hay que aclarar sus límites; se puede buscar financiamiento, pero no se puede sustituir el progreso real con demo irreproducibles. Para la inteligencia corporizada, la credibilidad de la industria es en sí misma un bien público.
\end{importantbox}

\begin{knowledgebox}{Matriz de validación: qué cuenta como progreso real}
\begin{tabularx}{\textwidth}{p{0.20\textwidth}p{0.34\textwidth}X}
\toprule
Elemento de validación & Qué hay que observar & Con qué no puede sustituirse \\
\midrule
Presentación pública & En vivo, sin teleoperación oculta, tareas reproducibles & Videos editados y fragmentos de un único intento exitoso. \\
Operación en escenarios & Cuántos pedidos o tareas reales se procesan al día y durante cuánto tiempo & Alianzas estratégicas y cartas de intención. \\
Escala de entregas & Si se alcanzan despliegues reales del orden de miles o decenas de miles de unidades & El número de prototipos de ingeniería. \\
Productividad & Producción por unidad de tiempo, tasa de fallas, proporción de respaldo humano & Apariencia y movimientos «parecidos a los humanos». \\
Volante de datos & Si las fallas del despliegue regresan al entrenamiento y a la mejora del producto & Centros de recolección estáticos y demostraciones únicas. \\
\bottomrule
\end{tabularx}
\end{knowledgebox}

Esta matriz también explica por qué Wang He insiste en no dañar la reputación de la industria. China enfrenta el envejecimiento de su población y un déficit de mano de obra; la productividad robótica no es solo una oportunidad de emprendimiento, sino que podría influir en la oferta a largo plazo de la manufactura y los servicios. Si la industria se ve arrastrada a una «era de hielo» por presentaciones falsas y promesas poco realistas, el problema social de quienes realmente necesitan robots para complementar la mano de obra también perderá una posible solución.

\subsection{El significado del episodio con Jensen Huang}

Lo anterior trató la credibilidad de la industria; esta sección usa el episodio final de la entrevista para volver al ecosistema tecnológico. Durante la visita de Jensen Huang (黄仁勋) a China, Wang He compartió mesa con él y se sentó a su lado; observó que Jensen Huang tolera el picante, que comió bastante carne de cerdo hervida en aceite picante y que respondió con entusiasmo a una función de cambio de máscaras. El fragmento parece anecdótico, pero en estas notas puede servir como indicio indirecto: NVIDIA, la IA física, la inteligencia corporizada y los emprendedores chinos de robótica ya se encuentran en el mismo campo industrial. El cómputo, la simulación, el cuerpo del robot y los escenarios de aplicación ya no son líneas paralelas.

La explicación que da Wang He de esa comida compartida también se relaciona con el eje técnico: NVIDIA concede mucha importancia a los datos sintéticos, porque estos requieren renderizado, simulación y entrenamiento en GPU; si esta ruta funciona, las empresas de GPU podrán sostener la mitad de la inteligencia corporizada. Galbot (银河通用) también mostró a NVIDIA su trabajo de entrenamiento con datos sintéticos y transferencia al mundo real. Esto muestra que «la productividad es el producto» y «los datos sintéticos son la clave de la recipe de datos» no son dos juicios separados, sino los dos extremos de un mismo ciclo cerrado: la producción de datos sostiene la capacidad, la capacidad entra en los escenarios y genera productividad, y los escenarios devuelven datos reales.

\subsection{Resumen de la sección}

El capital amplifica las oportunidades, pero también el desorden. Si la inteligencia corporizada quiere construir credibilidad a largo plazo, debe desplazar el reward del financiamiento y los demo hacia las presentaciones veraces, la productividad, la escala de entregas y el valor para el cliente.
